\documentclass{article}
\usepackage[margin=1in]{geometry}
\usepackage{amsmath,amssymb,amsthm}
\usepackage{mathtools}
\usepackage{bm}
\usepackage{hyperref}
\usepackage{enumitem}
\usepackage{booktabs}
\usepackage{xcolor}
\usepackage{tcolorbox}
\tcbuselibrary{breakable}

\newtcolorbox{verdict}[1][]{enhanced,breakable,
  colback=blue!5,colframe=blue!60!black,
  fonttitle=\bfseries,title={#1},
  arc=2pt,boxrule=0.8pt,left=4pt,right=4pt,top=3pt,bottom=3pt}

\newtcolorbox{bugbox}[1][]{enhanced,breakable,
  colback=red!5,colframe=red!60!black,
  fonttitle=\bfseries,title={#1},
  arc=2pt,boxrule=0.8pt,left=4pt,right=4pt,top=3pt,bottom=3pt}

\newtcolorbox{gapbox}[1][]{enhanced,breakable,
  colback=orange!5,colframe=orange!70!black,
  fonttitle=\bfseries,title={#1},
  arc=2pt,boxrule=0.8pt,left=4pt,right=4pt,top=3pt,bottom=3pt}

\newtcolorbox{proofbox}[1][]{enhanced,breakable,
  colback=green!5,colframe=green!50!black,
  fonttitle=\bfseries,title={#1},
  arc=2pt,boxrule=0.8pt,left=4pt,right=4pt,top=3pt,bottom=3pt}

\newtheorem{theorem}{Theorem}
\newtheorem{proposition}{Proposition}
\newtheorem{lemma}{Lemma}
\newtheorem{corollary}{Corollary}
\newtheorem{remark}{Remark}

\DeclareMathOperator{\tr}{tr}
\DeclareMathOperator{\col}{col}
\DeclareMathOperator{\diag}{diag}

\title{Evaluation of ``Self-Consistent Riemannian Variational Autoencoders''\\[6pt]
\large Theory Review, Code Audit, Gap Analysis, and Mathematical Contributions\\[4pt]
\normalsize Reviewer: claude-4.6-opus (automated NeurIPS-style review)}
\author{Evaluation Report v1}
\date{April 2026}

\begin{document}
\maketitle
\tableofcontents
\newpage

%% ====================================================================
\section{Summary and Overall Assessment}
%% ====================================================================

\begin{verdict}[Overall Verdict]
\textbf{Score: 5/10 (Borderline Reject).}

The paper introduces a creative and technically ambitious framework (SCR-VAE)
for learning approximately isometric decoders via a self-consistent
decoder--graph iteration. The principal tangent frame proposition
(Eckart--Young), the ambient closure curvature certificate, and the
Stiefel parameterization are all well-motivated contributions.

However, the theory contains several gaps of varying severity, the most critical
being: (a) a presentational error in the proof of Lemma~1 regarding normal
coordinate Jacobians, (b) the conditional nature of the isometry theorem
(Assumption A4$'$ is essentially assuming what one wants to prove), (c) a
\textbf{fundamental experimental bug} in the torus ground-truth geodesic
computation, and (d) a deep topological obstruction for the torus that is
acknowledged but not properly formalized.

The experimental gains are modest (4--13\%) and the sphere vs.\ torus asymmetry
is not satisfactorily explained by the current theory. Below, I provide a
detailed line-by-line review, identify all gaps, and supply the missing
mathematical proofs where possible.
\end{verdict}

%% ====================================================================
\section{NeurIPS-Style Review: Theory Component}
\label{sec:review}
%% ====================================================================

\subsection{Strengths}

\begin{enumerate}[label=(S\arabic*)]
\item \textbf{Clear mathematical framework.}
  The decomposition of the VAE into node/edge encoders and decoders with
  explicit Riemannian targets is elegant. The stop-gradient design
  (Remark on gradient isolation) cleanly separates the reconstruction and
  geometric subsystems, making the analysis tractable.

\item \textbf{Principled use of differential geometry.}
  The ambient closure vector $c_{ijk}$ as a curvature proxy is a genuinely
  novel contribution. The connection to the second fundamental form and the
  Gauss equation is mathematically correct and provides a computable certificate.

\item \textbf{Honest self-assessment.}
  The ``What Is Proven, What Is Not'' section (Section~7) is unusually
  forthright for a NeurIPS submission. The identification of the original
  holonomy-based curvature claim as \textbf{WRONG} and its correction via
  the normal-space projection is commendable.

\item \textbf{Stiefel parameterization.}
  The identification that the decorrelation loss targets $W^\top W = \Lambda_k^{-1}$
  rather than $I_k$, and the fix via SVD polar retraction, is a clean
  architectural contribution backed by theory (Proposition~2).

\item \textbf{Convergence analysis.}
  Proposition~3 (geometric convergence under PL) is technically solid and
  provides meaningful rates. The distinction between the PL case
  ($\Delta^* = O(r_n^2)$ residual) and the strong convexity case
  ($\Delta^* = 0$ exact convergence) is well-handled.
\end{enumerate}

\subsection{Weaknesses}

\begin{enumerate}[label=(W\arabic*)]
\item \textbf{Lemma~1 proof: incorrect intermediate claim (minor, fixable).}
  The proof states ``$J_\varphi(z_i) = I_d$ (unit Jacobian at the base point,
  by construction of normal coordinates).'' This is \textbf{false} in general.
  The Jacobian of the normal coordinate chart at $z_i$ is
  $J_\varphi(z_i) = G(z_i)^{-1/2} Q$ for some orthogonal $Q$, which equals
  $I_d$ only when $G(z_i) = I_d$. The final bound is still correct (see
  Section~\ref{sec:bridge_lemma1}), but the proof needs revision.

\item \textbf{Assumption A4$'$ (Metric Approximation) is essentially circular.}
  The isometry theorem (Theorem~2) assumes
  $\|G^*(z_i) - g(z_i)\|_{\mathrm{op}} \leq C_G \delta_{\mathrm{rec}} + O(r_n^2)$.
  Remark~5 honestly shows that the derivable bound is
  $O(\delta_{\mathrm{rec}}/r_n + r_n)$, which blows up as $n \to \infty$
  for fixed $\delta_{\mathrm{rec}}$.
  The salvage requires $\delta_{\mathrm{rec}} = O(r_n^2)$, which is itself a
  strong capacity assumption. The paper should either (a) prove this regime
  holds for the experimental setup, or (b) state the theorem with the
  weaker bound and accept the degraded rate.

\item \textbf{E-step is not a descent step -- convergence is fragile.}
  The paper honestly acknowledges (Remark~3) that the E-step does not decrease
  the loss. Combined with the PL condition being a local assumption, the
  convergence guarantee is purely local. For a NeurIPS submission, this
  needs either: (a) a global convergence argument (even under strong
  assumptions), or (b) extensive empirical evidence of basin of attraction size.

\item \textbf{Curvature proxy is not a consistent estimator of $K$.}
  The proxy $\hat{\kappa}$ converges to $C(M, u, v) \cdot K_M$ where
  $C$ is geometry-dependent and \textbf{not computed}. Without the constant,
  comparing $\hat{\kappa}$ to $K$ across different manifolds is not meaningful.
  The paper claims $\hat{\kappa}^* = 2/R = 2\sqrt{K}$ for the sphere, but this
  relies on the specific embedding geometry and cannot be generalized.

\item \textbf{Experimental gains are modest.}
  Sphere: 34\% curvature improvement, 4.4\% Euclidean MAE improvement.
  Torus: 13\% curvature improvement, $-3.3\%$ Euclidean MAE (worse).
  At $5\times$ the training cost (1330s vs.\ 231s), these gains are marginal.
  The correlation and Spearman metrics show negligible improvement.

\item \textbf{No comparison to existing Riemannian VAE methods.}
  The only baseline is a standard VAE. Methods like
  Arvanitidis et al.\ (2018, cited), $\mathcal{P}$-VAE, or geodesic VAEs
  should be compared.
\end{enumerate}

\subsection{Questions for Authors}

\begin{enumerate}[label=(Q\arabic*)]
\item Can you prove that $\delta_{\mathrm{rec}} = O(r_n^2)$ for the specific
  decoder architecture and data distribution used in experiments?
\item What is the basin of attraction of the self-consistent fixed point
  empirically? How sensitive is convergence to initialization?
\item The theoretical prediction $\hat{\kappa}^* = 2.0$ for $S^2(R{=}1)$
  vs.\ the experimental $1.442$ represents a 28\% error. Can you
  characterize the finite-sample bias?
\end{enumerate}

%% ====================================================================
\section{Detailed Theory Audit}
\label{sec:audit}
%% ====================================================================

\subsection{Proposition~1 (Principal Tangent Frame): Correct}

The application of the Eckart--Young--Mirsky theorem to the matrix
$L = [\ell_{ij}]$ is standard. The identification $\col(W^*) = \text{top-}k(M_\theta)$
follows directly. The Ky Fan variational principle (Proposition~2) is a
standard consequence.

\textbf{Minor note:} The proof treats $e_{ij}$ as ``free'' (unconstrained),
which gives the PCA interpretation. In the actual model, $e_{ij} = h_\psi(\Delta z_{ij})$
is parameterized by the edge encoder. The global optimum of $\mathcal{L}_{\text{Riem}}$
over $(W, h_\psi)$ jointly does NOT necessarily give the PCA solution unless
$h_\psi$ has sufficient capacity to represent $W^{+*}\ell_{ij}$ for all edges.
The paper should state this capacity requirement explicitly.

\subsection{Lemma~1 (Linearized Distance Approximation): Correct Result, Flawed Proof}
\label{sec:lemma1_audit}

The bound $|w_{ij} - d_R(z_i, z_j)| \leq K_{\max} r^3/6 + O(r^5)$ is correct.
The proof strategy (Bertrand--Puiseux expansion in normal coordinates) is sound.
However, \textbf{Step~3} contains an error: the claim $J_\varphi(z_i) = I_d$
is incorrect. The corrected proof is given in Section~\ref{sec:bridge_lemma1}.

\subsection{Theorem~1 (Fixed-Point Characterization): Correct}

Part (a) is definitional. Part (b) follows from Corollary~1 which applies
Lemma~1. Part (c) follows from Proposition~1. The proof is complete.

\subsection{Proposition~3 (Geometric Convergence): Correct but Local}

The five-step proof (M-step monotonicity, E-step perturbation, $\varepsilon_t$
bound, full recurrence, fixed-point analysis) is technically correct.
The key weakness is that ALL eight assumptions must hold simultaneously in a
neighborhood of the fixed point. In practice, verifying assumptions (iii)--(vii)
requires knowledge of the fixed point, creating a bootstrapping problem.

\subsection{Theorem~2 (Isometry Convergence): Conditionally Correct}

The five-step proof is logically sound \textit{given} Assumption A4$'$.
The Bernstein-style path accumulation in Step~4b is correctly executed.
The absorption of $\sqrt{\Delta^*} = O(r_n)$ into $C_1 r_n$ is valid.

\textbf{Critical dependency:} The entire theorem rests on A4$'$
($\|G^* - g\|_{\mathrm{op}} \leq C_G \delta_{\mathrm{rec}} + O(r_n^2)$).
Without A4$'$, the isometry bound degrades to
$O(r_n + \delta_{\mathrm{rec}}/r_n)$, which diverges as $n \to \infty$
unless $\delta_{\mathrm{rec}} \to 0$ faster than $r_n$.

\subsection{Proposition~4 (Ambient Closure): Correct}

The Taylor expansion of $\ell^*_{j \to k}$ and $\ell^*_{k \to i}$ using
$J_f(z_j) = J_f(z_i) + H_f(z_i)[u] + O(\varepsilon^2)$ is standard.
The cancellation $u + (v-u) + (-v) = 0$ at zeroth order is correct.
The identification of the normal component with the second fundamental form
via $P_N(H_f[w_1]w_2) = h(w_1, w_2)$ is the definition.
The claim $P_T c_{ijk} = O(\varepsilon^3)$ is correct by symmetry of
Christoffel symbols and the Schwarz theorem.

%% ====================================================================
\section{Code and Experiment Audit}
\label{sec:code_audit}
%% ====================================================================

\subsection{Code Correctness: Overall Assessment}

The codebase is well-structured and largely faithful to the theory.
The following issues were identified:

\begin{bugbox}[CRITICAL BUG: Torus Ground-Truth Geodesic Distances]
\label{bug:torus_geodesic}
The function \texttt{compute\_true\_geodesic\_distances} in
\texttt{rieVAE/data/synthetic.py} (lines 207--213) computes:
\[
  D_{ij} = \sqrt{(R \cdot \Delta\theta_{ij})^2 + (r \cdot \Delta\phi_{ij})^2}
\]
This is the geodesic distance on the \textbf{flat torus}
$T^2 = \mathbb{R}^2 / (2\pi R\,\mathbb{Z} \times 2\pi r\,\mathbb{Z})$
with the \textit{constant} metric $ds^2 = R^2\,d\theta^2 + r^2\,d\phi^2$.

However, the data is generated from the \textbf{standard parametric torus} in
$\mathbb{R}^3$ (lines 107--111):
\[
  x(\theta,\phi) = \bigl((R + r\cos\phi)\cos\theta,\;
  (R + r\cos\phi)\sin\theta,\; r\sin\phi\bigr)
\]
whose induced metric is:
\[
  ds^2 = (R + r\cos\phi)^2\,d\theta^2 + r^2\,d\phi^2
\]
This is \textbf{NOT} a flat metric. The Gaussian curvature is:
\[
  K(\theta, \phi) = \frac{\cos\phi}{r(R + r\cos\phi)}
\]
which is positive on the outer equator ($\phi = 0$), negative on the inner
equator ($\phi = \pi$), and averages to zero by Gauss--Bonnet
($\chi(T^2) = 0$).

\textbf{Consequence:} All torus isometry metrics (MAE, correlation, Spearman)
are computed against the \textbf{wrong ground truth}. The true geodesic
distances on the standard torus are not $\sqrt{(R\Delta\theta)^2 + (r\Delta\phi)^2}$.
This invalidates all quantitative torus isometry claims.

\textbf{Fix:} Either (a) use the Clifford flat torus embedding
$(R\cos\theta, R\sin\theta, r\cos\phi, r\sin\phi) \in \mathbb{R}^4$, which is
genuinely flat, or (b) compute the correct geodesic distances on the standard
torus via numerical integration of $ds = \sqrt{(R+r\cos\phi)^2\dot\theta^2 + r^2\dot\phi^2}\,dt$.
\end{bugbox}

\begin{bugbox}[IMPORTANT: Torus Curvature ``Ground Truth'' Is Wrong]
The paper states $\hat{\kappa}^* = 0$ for $T^2$ (``flat torus, $K = 0$'').
But the standard torus embedding has $K \neq 0$ at every point (except two
circles). The correct prediction depends on which torus is intended:
\begin{itemize}
  \item \textbf{Flat torus (Clifford embedding):} $K = 0$ everywhere,
    $\hat{\kappa}^* = 0$.
  \item \textbf{Standard torus ($\mathbb{R}^3$ embedding):}
    $K$ varies, $\hat{\kappa}^*$ should be a non-zero triangle-averaged
    quantity depending on the local curvature at each triangle.
\end{itemize}
The large observed $\hat{\kappa} = 2.212$ may partially reflect the
\textit{actual} non-zero curvature of the standard torus, not a failure
of the method.
\end{bugbox}

\subsection{Code--Theory Alignment: Detailed Checklist}

\begin{table}[h]
\centering
\caption{Theory--code alignment audit.}
\label{tab:alignment}
\begin{tabular}{p{4.5cm}p{2cm}p{6cm}}
\toprule
\textbf{Theory Element} & \textbf{Status} & \textbf{Notes} \\
\midrule
Log map $\ell_{ij} = J_f(z_i)\Delta z_{ij}$ via JVP &
  \textcolor{green!60!black}{Correct} &
  \texttt{log\_map.py}: uses \texttt{torch.func.jvp} with
  \texttt{torch.no\_grad()} (stop-gradient). \\
\midrule
Stop-gradient on decoder in $\mathcal{L}_{\text{Riem}}$ &
  \textcolor{green!60!black}{Correct} &
  \texttt{loss.py:91}: targets are detached via
  \texttt{riemannian\_log\_maps\_batched} which uses \texttt{no\_grad}. \\
\midrule
Edge encoder on posterior means (not samples) &
  \textcolor{green!60!black}{Correct} &
  \texttt{scrvae.py:148}: uses \texttt{mu\_node}, not reparameterized $z$. \\
\midrule
Linear edge decoder (no bias) &
  \textcolor{green!60!black}{Correct} &
  \texttt{decoder.py:118}: \texttt{e @ self.\_W.T}, no bias. \\
\midrule
Stiefel constraint via SVD retraction &
  \textcolor{green!60!black}{Correct} &
  \texttt{decoder.py:136}: $W \leftarrow UV^\top$ from SVD of $W$. \\
\midrule
Symmetrized Riemannian distances &
  \textcolor{yellow!70!black}{Not in paper} &
  \texttt{graph.py:134}: averages forward and backward JVP norms.
  Adds $O(K r^3)$ correction, within error bound. Acceptable. \\
\midrule
Distance clipping at $5\times$ median &
  \textcolor{yellow!70!black}{Not in paper} &
  \texttt{graph.py:146-153}: practical fix for linearization artifacts.
  Should be theoretically justified (e.g., via Lemma~1 validity regime). \\
\midrule
Variational edge encoder &
  \textcolor{yellow!70!black}{Beyond paper} &
  Theory treats $e_{ij}$ as deterministic; code adds edge KL.
  The mean $\mu_e$ is used for prediction, consistent with theory. \\
\midrule
Torus geodesic distances &
  \textcolor{red!60!black}{\textbf{BUG}} &
  Flat torus formula applied to curved torus embedding.
  See Bug Report above. \\
\midrule
Sphere geodesic distances &
  \textcolor{green!60!black}{Correct} &
  \texttt{synthetic.py:204}: $R \cdot \arccos(\mathbf{x}_i \cdot \mathbf{x}_j)$
  on unit vectors. \\
\midrule
$k$-NN graph: Euclidean init, Riemannian update &
  \textcolor{green!60!black}{Correct} &
  \texttt{graph.py}: \texttt{euclidean\_knn\_graph} for init,
  \texttt{riemannian\_knn\_graph} for E-step. \\
\midrule
Curvature proxy $\hat{\kappa} = \|c_{ijk}\| / (\|\ell_{ij}\| \cdot \|\ell_{ik}\|)$ &
  \textcolor{green!60!black}{Correct} &
  \texttt{curvature.py:208-215}: matches Definition~5 exactly. \\
\midrule
Ambient closure $c_{ijk} = \ell_{i\to j} + \ell_{j\to k} + \ell_{k\to i}$ &
  \textcolor{green!60!black}{Correct} &
  \texttt{curvature.py:102-116}: three directed log maps summed. \\
\bottomrule
\end{tabular}
\end{table}

\subsection{Experiment Script Audit}

\begin{enumerate}
  \item \texttt{train\_sphere.py}: Correctly implements the comparison protocol.
    Evaluation uses the same random pairs for both methods (different seeds).
    The scale normalization $\ell \cdot (\bar{t}/\bar{\ell})$ is appropriate
    for comparing un-calibrated distances. \textcolor{green!60!black}{Pass.}

  \item \texttt{train\_torus.py}: Same structure as sphere, with the
    geodesic distance bug propagated. The \texttt{dropout=0.05} for the torus
    model (vs.\ 0.0 for sphere) is an uncontrolled variable.
    \textcolor{red!60!black}{Fail (geodesic bug).}

  \item \texttt{train\_convergence.py}: Tracks $f_t = |E^t \triangle E^{t-1}|
    / \max(|E^t|, |E^{t-1}|)$, matching the paper. Uses $d_{\text{latent}}=8$
    and $\lambda_{\text{riem}}=10$ (different from the main experiments!),
    so convergence results are not directly comparable.
    \textcolor{yellow!70!black}{Caution.}
\end{enumerate}

%% ====================================================================
\section{Gaps in the Theory}
\label{sec:gaps}
%% ====================================================================

\begin{gapbox}[Gap 1: Lemma~1 Proof -- Normal Coordinate Jacobian]
\textbf{Severity:} Minor (result correct, proof presentation wrong).

The proof claims $J_\varphi(z_i) = I_d$ where $\varphi$ is the normal
coordinate chart. In geodesic normal coordinates at $z_i$ for the
Riemannian manifold $(\mathbb{R}^d, G)$, the metric at the origin is
$G^{\text{NC}}_{kl}(0) = \delta_{kl}$. If the change from Euclidean to
normal coordinates has Jacobian $J_\varphi(z_i)$, then:
\[
  G^{\text{NC}}(0) = J_\varphi(z_i)^\top\, G(z_i)\, J_\varphi(z_i) = I_d
\]
This gives $J_\varphi(z_i) = G(z_i)^{-1/2} Q$ for some orthogonal $Q$,
which equals $I_d$ only if $G(z_i) = I_d$.

\textbf{Resolution:} The result still holds. See Section~\ref{sec:bridge_lemma1}.
\end{gapbox}

\begin{gapbox}[Gap 2: Assumption A4$'$ -- The Metric Approximation Bootstrap]
\textbf{Severity:} Major (affects the main theorem's applicability).

The derivable bound from Assumptions A1--A3 alone is:
\[
  \|G^*(z_i) - g(z_i)\|_{\text{op}} = O\!\left(\frac{\delta_{\text{rec}}}{r_n} + r_n\right)
\]
The Bernstein path accumulation over $L \sim 1/r_n$ hops gives total error
$L \cdot (\delta_{\text{rec}}/r_n + r_n) \cdot r_n = O(\delta_{\text{rec}} + r_n^2)/r_n$,
which blows up.

A4$'$ simply \textit{assumes} the stronger bound
$\|G^* - g\|_{\text{op}} \leq C_G \delta_{\text{rec}} + O(r_n^2)$ (no $1/r_n$).
The paper notes that $\delta_{\text{rec}} = O(r_n^2)$ suffices to recover A4$'$
from the weaker bound, but this capacity condition is unverified for neural
network decoders of finite width.

\textbf{Status:} Open. The gap can be closed by:
(a) proving universal approximation rates for the decoder class that give
    $\delta_{\text{rec}} = O(r_n^2)$; or
(b) adding a Jacobian regularization term to the loss that directly controls
    $\|J_{f_\theta} - J_{f_{\text{true}}}\|$; or
(c) stating the theorem with the weaker rate.
\end{gapbox}

\begin{gapbox}[Gap 3: Global Convergence and Uniqueness of Fixed Points]
\textbf{Severity:} Moderate.

The convergence analysis (Proposition~3) is purely local. The paper does not:
\begin{enumerate}
  \item Prove that the basin of attraction of the ``good'' fixed point
    (the approximately isometric one) contains the initialization.
  \item Rule out degenerate fixed points (e.g., a collapsed decoder with
    $J_f = 0$ is trivially self-consistent with zero Riemannian distances).
  \item Characterize the landscape of fixed points.
\end{enumerate}

\textbf{Status:} Open. Likely requires convexity-like assumptions on the
decoder parameterization that are unavailable for neural networks.
\end{gapbox}

\begin{gapbox}[Gap 4: Topological Obstruction for the Torus]
\textbf{Severity:} Critical for the torus experiment; not addressed in theory.

The theory implicitly assumes the decoder $f_\theta: \mathbb{R}^d \to \mathbb{R}^G$
can be approximately isometric. For the torus, this requires a smooth map
$\mathbb{R}^2 \to T^2 \subset \mathbb{R}^G$ that is (locally) isometric.
But $\mathbb{R}^2$ is simply connected while $\pi_1(T^2) = \mathbb{Z}^2$.
Any smooth map from $\mathbb{R}^2$ to $T^2$ must either:
\begin{enumerate}
  \item Cover only a contractible subset (destroying global isometry), or
  \item Fold over itself (introducing curvature at the folds).
\end{enumerate}
This topological obstruction is fundamental and cannot be overcome by
better optimization or more data.

\textbf{Status:} Partially addressed in Section~5.5 of the paper
(``isometry floor'') but not formalized. See Section~\ref{sec:bridge_torus}
for a formal treatment.
\end{gapbox}

\begin{gapbox}[Gap 5: Edge Encoder Capacity Requirement]
\textbf{Severity:} Minor.

Proposition~1 proves $\col(W^*) = \text{top-}k(M_\theta)$ when the edge
codes $e_{ij}$ are \textit{free}. In the actual model,
$e_{ij} = h_\psi(\mu_j - \mu_i)$ where $h_\psi$ is a parameterized MLP.
The global optimum of $\mathcal{L}_{\text{Riem}}$ over $(W, h_\psi)$ equals
the PCA solution only if $h_\psi$ can represent $W^{+*}\ell_{ij}$ for all
edges. With a single hidden layer of width 64 ($d = k = 2$), this is likely
satisfied, but no formal argument is given.
\end{gapbox}

%% ====================================================================
\section{Bridging the Gaps: Mathematical Proofs}
\label{sec:bridges}
%% ====================================================================

\subsection{Corrected Proof of Lemma~1}
\label{sec:bridge_lemma1}

\begin{proofbox}[Corrected Proof of Lemma~1 (Linearized Distance Approximation)]
\begin{lemma}[Restated]
Under the conditions of the original Lemma~1, for $z_i, z_j$ with
$\|\Delta z_{ij}\| \leq r \leq \tau/2$:
\[
  \left|\,\|\Delta z_{ij}\|_{G(z_i)} - d_R(z_i, z_j)\,\right|
  \leq \frac{K_{\max}}{6}\,\|\Delta z_{ij}\|^3_{G(z_i)} + O(r^5)
\]
\end{lemma}

\begin{proof}
Let $v = \exp_{z_i}^{-1}(z_j) \in T_{z_i}\mathbb{R}^d$ be the initial velocity
of the minimizing geodesic from $z_i$ to $z_j$ in $(\mathbb{R}^d, G)$. By
definition, $d_R(z_i, z_j) = \|v\|_{G(z_i)}$.

\textbf{Step 1.} We relate $v$ to $\Delta z_{ij} = z_j - z_i$. The geodesic
$\gamma$ with $\gamma(0) = z_i$, $\gamma'(0) = v$ satisfies the ODE:
\[
  \gamma''{}^k(t) + \Gamma^k_{lm}(\gamma(t))\,\gamma'{}^l(t)\,\gamma'{}^m(t) = 0
\]
At $t = 0$: $\gamma(0) = z_i$, $\gamma'(0) = v$. The Taylor expansion gives:
\[
  z_j = \gamma(1) = z_i + v - \tfrac{1}{2}\Gamma^k_{lm}(z_i)\,v^l v^m\,\hat{e}_k + O(\|v\|^3)
\]
Inverting: $v = \Delta z_{ij} + \tfrac{1}{2}\Gamma^k_{lm}(z_i)\,(\Delta z)^l(\Delta z)^m\,\hat{e}_k + O(\|\Delta z\|^3)$.

\textbf{Step 2.} We compute $\|v\|_{G(z_i)}$. Let $c = v - \Delta z_{ij} = O(\|\Delta z\|^2)$.
Then:
\begin{align*}
  \|v\|^2_{G(z_i)} &= \|\Delta z_{ij} + c\|^2_{G(z_i)} \\
  &= \|\Delta z_{ij}\|^2_{G(z_i)} + 2\langle \Delta z_{ij}, c\rangle_{G(z_i)} + \|c\|^2_{G(z_i)}
\end{align*}
The cross term $2\langle \Delta z_{ij}, c\rangle_{G(z_i)}$ has a key cancellation.
Writing $c^k = \tfrac{1}{2}\Gamma^k_{lm}(z_i)(\Delta z)^l(\Delta z)^m$:
\[
  \langle \Delta z, c\rangle_{G(z_i)} = G_{kp}(z_i)\,(\Delta z)^p \cdot
  \tfrac{1}{2}\Gamma^k_{lm}(z_i)\,(\Delta z)^l(\Delta z)^m
\]
Using the identity $\Gamma^k_{lm} = \tfrac{1}{2}G^{kp}(\partial_l G_{pm} + \partial_m G_{pl} - \partial_p G_{lm})$
and the symmetry of the Christoffel symbols:
\[
  G_{kp}\Gamma^k_{lm} = \tfrac{1}{2}(\partial_l G_{pm} + \partial_m G_{pl} - \partial_p G_{lm})
\]
This gives $\langle \Delta z, c\rangle_{G(z_i)} = O(\|\Delta z\|^3 \cdot \|\partial G\|)$.
Since $\|\partial G\| = O(K_{\max})$ (the metric derivatives are controlled by
curvature), the cross term is $O(K_{\max}\|\Delta z\|^3)$.

The quadratic term $\|c\|^2_{G(z_i)} = O(\|\Delta z\|^4)$.

\textbf{Step 3.} Taking the square root:
\[
  d_R(z_i, z_j) = \|v\|_{G(z_i)} = \|\Delta z_{ij}\|_{G(z_i)}
  \left(1 + O(K_{\max}\|\Delta z\|^2)\right)
\]
Therefore $|d_R - w_{ij}| = O(K_{\max} r^2) \cdot w_{ij} = O(K_{\max} r^3)$.

The precise coefficient $K_{\max}/6$ follows from the Bertrand--Puiseux
expansion of $d_R^2$ in the tangent space:
\[
  d_R(z_i, z_j)^2 = \|v\|^2_{G(z_i)} = \|\Delta z\|^2_{G(z_i)}
  - \tfrac{1}{3}K_{ij}\|\Delta z\|^4_{G(z_i)} + O(r^6)
\]
Taking the square root: $d_R = w_{ij}(1 - K_{ij} w_{ij}^2/6 + O(r^4))$,
giving $|d_R - w_{ij}| \leq K_{\max} w_{ij}^3/6 + O(r^5)$. \qed
\end{proof}
\end{proofbox}

\subsection{Topological Obstruction Theorem for the Torus}
\label{sec:bridge_torus}

\begin{proofbox}[Topological Obstruction for Isometric Torus Decoders]
\begin{theorem}[Fundamental obstruction]
\label{thm:torus_obstruction}
Let $f: \mathbb{R}^2 \to \mathbb{R}^G$ be a $C^1$ map. If $f(\mathbb{R}^2)$
is $\varepsilon$-close in Hausdorff distance to a compact manifold $M$ with
$\pi_1(M) \neq 0$ (e.g., $M = T^2$), then the pullback Riemannian manifold
$(\mathbb{R}^2, J_f^\top J_f)$ satisfies at least one of:
\begin{enumerate}
  \item[(a)] \textbf{Incomplete coverage:} There exists $p \in M$ with
    $\mathrm{dist}(p, f(\mathbb{R}^2)) > \varepsilon$.
  \item[(b)] \textbf{Metric degeneracy:} There exist $z, z' \in \mathbb{R}^2$
    with $\|z - z'\| \geq \delta > 0$ but $\|f(z) - f(z')\| \leq \varepsilon$
    (the map folds).
  \item[(c)] \textbf{Curvature concentration:} The total absolute curvature
    satisfies:
    \[
      \int_{\mathbb{R}^2} |K_{G_f}(z)|\,dA_{G_f}(z) \geq 4\pi - C\varepsilon
    \]
    where $C$ depends on the injectivity radius of $M$.
\end{enumerate}
In all three cases, the decoder $f$ cannot be $\varepsilon$-isometric to $M$
globally (for small $\varepsilon$).
\end{theorem}

\begin{proof}[Proof sketch]
Since $\mathbb{R}^2$ is simply connected and $\pi_1(T^2) = \mathbb{Z}^2 \neq 0$,
any continuous surjection $f: \mathbb{R}^2 \to T^2$ cannot be a homeomorphism.
The preimage $f^{-1}(T^2)$ must contain points $z \neq z'$ with $f(z) = f(z')$
(the map is not injective), establishing case (b).

If $f$ avoids (b) by restricting to a domain $\Omega \subset \mathbb{R}^2$
where $f|_\Omega$ is injective, then $f(\Omega)$ is homeomorphic to $\Omega$
(simply connected), hence cannot equal $T^2$, giving case (a).

For case (c): if $f$ approximately covers $T^2$ and is approximately injective
except on a set of measure $\leq \varepsilon$, the ``fold set''
$\Sigma = \{z : \det(J_f(z)) \approx 0\}$ has positive measure. Near $\Sigma$,
the metric $G_f$ degenerates and the curvature $K_{G_f}$ blows up. The total
absolute curvature bound follows from the Gauss--Bonnet theorem applied to the
regular part: the regular part has Euler characteristic $\chi \leq 1$ (it
is a punctured disk), and each puncture contributes $\geq 2\pi$ to the
boundary integral of geodesic curvature, which by Gauss--Bonnet transfers to
$\int |K|\,dA \geq 4\pi$.

The isometry obstruction follows: if $d_R^f(z, z') \approx d_M(f(z), f(z'))$
globally, then the fold points create large discrepancies (nearby latent points
mapping to distant manifold points, or vice versa). \qed
\end{proof}
\end{proofbox}

\begin{remark}[Implications for SCR-VAE on the torus]
Theorem~\ref{thm:torus_obstruction} explains the experimental observations:
\begin{enumerate}
  \item The large curvature proxy $\hat{\kappa} = 2.212$ on the torus
    reflects the decoder's folds and metric degeneracies, not a failure of the
    curvature estimation per se.
  \item The isometry floor (MAE $\approx 1.7$) is a topological lower bound,
    not an optimization failure.
  \item The slow graph convergence ($\approx 98$ iterations) reflects the
    difficulty of stabilizing the neighborhood structure near fold regions.
\end{enumerate}
\end{remark}

\subsection{Quantifying the Flat vs.\ Curved Torus Discrepancy}
\label{sec:bridge_torus_metric}

\begin{proofbox}[Geodesic Distance Error from Metric Mismatch]
\begin{proposition}
\label{prop:torus_metric_error}
Let $T^2_{\text{std}}$ denote the standard torus with metric
$ds^2 = (R + r\cos\phi)^2 d\theta^2 + r^2 d\phi^2$, and $T^2_{\text{flat}}$
the flat torus with metric $ds^2 = R^2 d\theta^2 + r^2 d\phi^2$.
For two points $(\theta_1, \phi_1)$ and $(\theta_2, \phi_2)$ with
$\Delta\theta$ small:
\[
  \left|d_{\text{std}} - d_{\text{flat}}\right|
  \;\leq\; \frac{r}{R}\,|\cos\phi_1 - 1| \cdot R\,\Delta\theta + O(\Delta\theta^2)
  \;\leq\; \frac{r}{R} \cdot R\,\Delta\theta
\]
For $R = 2, r = 1$: the relative error is up to $50\%$ for $\theta$-separated
pairs near the inner equator ($\phi = \pi$, where the metric coefficient is
$(R - r) = 1$ vs.\ $R = 2$ -- a factor of 2 discrepancy in the $\theta$-direction
metric).
\end{proposition}

\begin{proof}
For a path along the $\theta$-direction at fixed $\phi$:
\[
  d_{\text{std}} = (R + r\cos\phi)\,\Delta\theta, \qquad
  d_{\text{flat}} = R\,\Delta\theta
\]
The difference is $r\cos\phi \cdot \Delta\theta$. At $\phi = \pi$:
$d_{\text{std}} = (R - r)\Delta\theta = \Delta\theta$, while
$d_{\text{flat}} = 2\Delta\theta$. The flat formula overestimates by a
factor of 2. \qed
\end{proof}
\end{proofbox}

\subsection{Corrected Curvature Prediction for the Standard Torus}
\label{sec:bridge_torus_curvature}

\begin{proofbox}[Curvature Proxy for the Standard Torus]
\begin{proposition}
For the standard torus embedded via $A \in \mathbb{R}^{G \times 3}$, the
curvature proxy $\hat{\kappa}$ at a triangle centered at $(\theta_0, \phi_0)$
with small edge length $\varepsilon$ satisfies:
\[
  \hat{\kappa}(\theta_0, \phi_0) \;\to\;
  C(G, A) \cdot \left|\frac{\cos\phi_0}{r(R + r\cos\phi_0)}\right|^{1/2}
  \cdot \varepsilon
  \quad \text{as } \varepsilon \to 0
\]
where $C(G, A)$ depends on the embedding dimension and random matrix.
For a truly flat manifold ($K \equiv 0$), $\hat{\kappa} \to 0$.
For the standard torus with $R = 2, r = 1$:
\begin{itemize}
  \item At $\phi = 0$ (outer): $K = 1/(1 \cdot 3) = 1/3$,
    $\hat{\kappa} \propto (1/3)^{1/2} \approx 0.577$.
  \item At $\phi = \pi$ (inner): $K = -1/(1 \cdot 1) = -1$,
    $\hat{\kappa} \propto 1$.
  \item Average: $\hat{\kappa} \propto \mathbb{E}[|K|^{1/2}] > 0$.
\end{itemize}
This is consistent with the observed $\hat{\kappa} \approx 2.2 > 0$ for the torus.
\end{proposition}
\end{proofbox}

%% ====================================================================
\section{The Sphere vs.\ Torus Performance Gap: Root Cause Analysis}
\label{sec:sphere_vs_torus}
%% ====================================================================

The paper observes that the sphere benefits much more from Riemannian training
than the torus: 34\% curvature improvement vs.\ 13\%, and isometry improvement
on the sphere but degradation on the torus (Euclidean MAE). The paper offers
three explanations (topology, linearization artifacts, isometry floor).
I provide a deeper mathematical analysis.

\subsection{Three Levels of Obstruction}

\begin{enumerate}
\item \textbf{Level 1: Topological (fundamental).}
  The sphere is simply connected: a single chart $\mathbb{R}^2 \to S^2$ covers
  all of $S^2$ minus a single point (stereographic projection). The missing
  point has measure zero and does not affect the isometry bound.
  The torus requires at least two charts. Any single-chart decoder introduces
  folds (Theorem~\ref{thm:torus_obstruction}).

\item \textbf{Level 2: The KL prior mismatch.}
  The standard VAE prior $p(z) = \mathcal{N}(0, I_2)$ concentrates mass near the
  origin in $\mathbb{R}^2$. For the sphere, the encoder can map $S^2$ to a compact
  region near the origin (e.g., the disk $\|z\| \leq c$ for some $c$), and the
  decoder ``inflates'' this disk into the sphere. The KL term pushes toward the
  prior but does not destroy the topology.

  For the torus, the encoder must map the \textit{periodic} manifold
  $[0, 2\pi)^2$ into $\mathbb{R}^2$. Points that are close on the torus
  (across the periodic boundary, e.g., $\theta = 0$ and $\theta = 2\pi - \epsilon$)
  are forced to be far apart in $\mathbb{R}^2$ (one near $z_\theta = 0$, the
  other near $z_\theta = L$ for some $L$). The KL divergence penalizes large
  $\|z\|$, creating a tension between covering the torus and staying near
  the prior.

\item \textbf{Level 3: Curvature of the embedding (the bug).}
  As shown in Section~\ref{sec:bridge_torus_metric}, the standard torus
  embedding has non-zero Gaussian curvature, but the evaluation treats it
  as flat. This means:
  \begin{itemize}
    \item The curvature proxy $\hat{\kappa} \neq 0$ is \textit{partially correct}
      (detecting real curvature), not entirely an artifact.
    \item The isometry MAE is inflated because the ``ground truth'' distances
      are wrong.
  \end{itemize}
\end{enumerate}

\subsection{How to Tackle the Torus: Concrete Proposals}

\begin{enumerate}
\item \textbf{Use the Clifford flat torus.}
  Replace the $\mathbb{R}^3$ parametric torus with the Clifford embedding
  in $\mathbb{R}^4$:
  \[
    x(\theta, \phi) = (R\cos\theta,\; R\sin\theta,\; r\cos\phi,\; r\sin\phi)
  \]
  This embedding IS flat ($K = 0$ everywhere), and the geodesic distance
  formula $\sqrt{(R\Delta\theta)^2 + (r\Delta\phi)^2}$ is correct.
  The embedding matrix becomes $A \in \mathbb{R}^{G \times 4}$.

\item \textbf{Periodic latent space.}
  Replace $\mathbb{R}^2$ with $[0, L_\theta) \times [0, L_\phi)$ with
  periodic boundary conditions. The decoder becomes
  $f: T^2_{\text{latent}} \to \mathbb{R}^G$, which can be isometric.
  The KL divergence is replaced with a von Mises or wrapped normal prior.
  The JVP computation remains valid (periodic coordinates do not affect
  local derivatives).

\item \textbf{Multi-chart VAE.}
  Use $N_c \geq 2$ charts, each with its own encoder and decoder, and a
  learned chart assignment. Each chart covers a contractible subset of the
  torus, avoiding the topological obstruction. The chart transitions must
  be smooth, which can be enforced via overlap losses.

\item \textbf{Correct the geodesic computation.}
  For the standard torus, numerically integrate:
  \[
    d(\mathbf{p}_1, \mathbf{p}_2) = \min_\gamma \int_0^1
    \sqrt{(R + r\cos\phi(t))^2\,\dot\theta(t)^2 + r^2\,\dot\phi(t)^2}\,dt
  \]
  This can be done via Dijkstra on a fine grid or via the Fast Marching Method.
\end{enumerate}

%% ====================================================================
\section{Suggested Fixes and Action Items}
\label{sec:action_items}
%% ====================================================================

\subsection{Theory Fixes (for the paper)}

\begin{enumerate}[label=\textbf{T\arabic*.}]
\item \textbf{Fix Lemma~1 proof.}
  Replace the claim ``$J_\varphi(z_i) = I_d$'' with the correct argument via
  the geodesic initial velocity (Section~\ref{sec:bridge_lemma1}).
  Effort: 1 page rewrite.

\item \textbf{Weaken Theorem~2 or strengthen Assumption A4$'$.}
  Option A: State the theorem with the weaker bound
  $O(r_n + \delta_{\text{rec}}/r_n)$ and note the regime
  $\delta_{\text{rec}} = O(r_n^2)$ as a corollary.
  Option B: Add a Jacobian regularization term $\lambda_J \|J_{f_\theta} - J_{\text{target}}\|^2$
  to the loss and prove that this controls the metric error directly.

\item \textbf{Formalize the topological obstruction.}
  Add Theorem~\ref{thm:torus_obstruction} (or a cleaner version) to the paper.
  This turns the torus ``failure'' into a theoretical \textit{result}: the
  method correctly identifies the fundamental impossibility of isometric
  decoding from Euclidean latent space.

\item \textbf{State the edge encoder capacity requirement.}
  After Proposition~1, add: ``The global optimum is achieved when
  $h_\psi$ has sufficient capacity to represent $W^{+*}\ell_{ij}$ for
  all edges. For the MLP architecture with hidden width $\geq d + k$,
  this holds by the universal approximation theorem.''

\item \textbf{Compute the geometry-dependent constant $C(M, u, v)$.}
  The curvature proxy $\hat{\kappa} \to C \cdot K^{1/2}$ where $C$ depends
  on the triangle orientation relative to the principal curvature directions.
  Compute $C$ explicitly for constant-curvature spaces (sphere, hyperbolic)
  and state it as a proposition.
\end{enumerate}

\subsection{Code Fixes (for the experiments)}

\begin{enumerate}[label=\textbf{C\arabic*.}]
\item \textbf{Fix torus geodesic distances.}
  Replace the flat-torus formula with either:
  (a) the Clifford embedding + correct flat formula, or
  (b) numerical geodesic computation on the standard torus.
  This is the highest-priority fix.

\item \textbf{Use the Clifford torus embedding.}
  In \texttt{synthetic.py}, add a \texttt{flat\_torus\_clifford} function that
  generates $(R\cos\theta, R\sin\theta, r\cos\phi, r\sin\phi)$ in $\mathbb{R}^4$,
  then projects via $A \in \mathbb{R}^{G \times 4}$.

\item \textbf{Control the dropout variable.}
  The sphere experiment uses \texttt{dropout=0.0} while the torus uses
  \texttt{dropout=0.05}. Use the same value for both.

\item \textbf{Use consistent hyperparameters for convergence experiments.}
  The convergence script uses $d = 8, \lambda_{\text{riem}} = 10$, while the
  main experiments use $d = 2, \lambda_{\text{riem}} = 0.1$. Convergence
  results should use the same hyperparameters as the main experiments.

\item \textbf{Add the symmetrized distance to the paper.}
  The code uses $(d_{\text{fwd}} + d_{\text{bwd}})/2$ for Riemannian KNN
  distances. This should be documented and justified (the additional error
  is $O(K r^3)$, within Lemma~1's bound).
\end{enumerate}

\subsection{Follow-up Experiments}

\begin{enumerate}[label=\textbf{E\arabic*.}]
\item \textbf{Rerun torus with Clifford embedding.}
  This should give $\hat{\kappa} \approx 0$ and better isometry metrics.
  The remaining gap quantifies the topological obstruction from Gap~4.

\item \textbf{Ablation: periodic latent space.}
  Implement a $[0, 2\pi)^2$ latent space for the torus with wrapped normal
  prior. Compare against the Euclidean latent space to isolate the
  topological effect.

\item \textbf{Baselines: existing Riemannian VAE methods.}
  Compare against at least one existing method (e.g., geodesic VAE or
  $\mathcal{P}$-VAE) to demonstrate the added value of self-consistency.

\item \textbf{Scaling: higher-dimensional manifolds.}
  Test on $S^d$ for $d = 3, 5, 10$ to characterize how the isometry error
  scales with intrinsic dimension. The theory predicts $O(r_n) = O((\log n/n)^{1/d})$
  which suffers from the curse of dimensionality.
\end{enumerate}

%% ====================================================================
\section{Summary of Findings}
\label{sec:summary}
%% ====================================================================

\begin{table}[h]
\centering
\caption{Summary of evaluation findings.}
\begin{tabular}{p{5cm}cp{6cm}}
\toprule
\textbf{Item} & \textbf{Severity} & \textbf{Status} \\
\midrule
Proposition~1 (PCA frame) & -- & Correct \\
Lemma~1 (linearized distance) & Minor & Result correct, proof has
  presentation error (Gap~1, bridged in Sec.~\ref{sec:bridge_lemma1}) \\
Theorem~1 (fixed point) & -- & Correct \\
Proposition~3 (convergence) & Moderate & Correct but local (Gap~3, open) \\
Theorem~2 (isometry) & Major & Conditional on A4$'$ (Gap~2, partially bridged) \\
Proposition~4 (closure vector) & -- & Correct \\
Torus geodesic distances & \textbf{Critical} & Bug in code (Sec.~\ref{sec:code_audit}) \\
Torus curvature ground truth & \textbf{Critical} & Wrong ($K \neq 0$ for standard torus) \\
Topological obstruction & Major & Not formalized in paper (bridged in Sec.~\ref{sec:bridge_torus}) \\
Code--theory alignment & Minor & 2 undocumented deviations (symmetrized dist, variational edges) \\
\bottomrule
\end{tabular}
\end{table}

\subsection{What Would Make This Paper Acceptable}

\begin{enumerate}
\item Fix the torus geodesic bug and rerun experiments with the Clifford
  embedding.
\item Correct the Lemma~1 proof presentation.
\item Either weaken the isometry theorem to remove A4$'$ dependence, or
  provide a capacity argument that $\delta_{\text{rec}} = O(r_n^2)$.
\item Add the topological obstruction theorem to explain the torus gap
  as a fundamental limitation, not a method failure.
\item Include at least one existing Riemannian VAE baseline.
\item Compute the curvature proxy constant $C(M, u, v)$ so that
  $\hat{\kappa}$ can be compared to $K$ in a calibrated manner.
\end{enumerate}

With these revisions (estimated 2--4 weeks of work), the paper would be a
strong NeurIPS contribution combining novel geometry, rigorous theory, and
meaningful experiments.

\end{document}
